\subsection{Self-Attention for Temporal Representations}

Self-attention mechanisms have become a standard tool for capturing long-range dependencies in sequential data. In a typical formulation, a sequence of representations $\mathbf{h} = [h_1, h_2, \ldots, h_T]$ is processed by computing relevance scores between pairs of positions, e.g., scaled dot-product similarity, and then forming context vectors as weighted sums over the sequence. Multi-head self-attention extends this by applying several such operations in parallel with different linear projections, allowing the model to attend to different aspects of the sequence.

Related work on combining LSTM with attention for time series includes the model of~\cite{AIKP2023LSTMmultiattn}, which advances multivariate forecasting by using an LSTM encoder followed by \emph{multiple} attention mechanisms. The LSTM first produces a sequence of hidden states $h_1, h_2, \ldots, h_T$ that encode temporal structure (trends, seasonality, lag relationships). Attention is then applied in two ways: \emph{temporal} attention, which weights which past timesteps matter most for prediction, and \emph{feature} attention, which weights which input variables contribute most at each time. Thus the model learns a joint importance over time steps and variables rather than treating all inputs equally, improving forecast accuracy. In L-GTA we also combine an LSTM encoder with attention: we apply \emph{temporal} self-attention over the sequence of hidden states and, optionally, \emph{channel} (cross-series) self-attention so that each series can attend to others; our goal is per-timestep latent generation and transformation rather than forecasting.

In L-GTA, we use a \emph{temporal multi-head self-attention} block over the Bi-LSTM encoder hidden states. The attention layer is applied in a deterministic way: it enriches the hidden states, and the resulting representation is then mapped by $\phi$ to the parameters $(\mu_t, \Sigma_t)$ of the CVAE latent distribution, from which $v_t$ is sampled. Thus, the stochasticity in our model is confined to the latent variable $v_t$ (the CVAE formulation); the self-attention mechanism itself is deterministic. This differs from approaches such as the Variational Self-Attention Mechanism (VSAM)~\cite{8614232}, which model the context vectors produced by attention as random variables with a learned posterior and integrate the reparameterization trick at the attention level. We instead keep standard self-attention and rely on the CVAE to provide a latent space suitable for transformation and generation.